\documentclass[11pt]{article}


\usepackage{fullpage}
\usepackage{graphicx}
\usepackage{amsmath}
\usepackage{amssymb}
\usepackage{amsthm}
\usepackage{fancyvrb}

\newcommand{\myname}{Mehshan Mustafa}

\newenvironment{theorem}[2][Theorem]{\begin{trivlist}
\item[\hskip \labelsep {\bfseries #1}\hskip \labelsep {\bfseries #2.}]}{\end{trivlist}}
\newenvironment{lemma}[2][Lemma]{\begin{trivlist}
\item[\hskip \labelsep {\bfseries #1}\hskip \labelsep {\bfseries #2.}]}{\end{trivlist}}
\newenvironment{exercise}[2][Exercise]{\begin{trivlist}
\item[\hskip \labelsep {\bfseries #1}\hskip \labelsep {\bfseries #2.}]}{\end{trivlist}}
\newenvironment{problem}[2][Problem]{\begin{trivlist}
\item[\hskip \labelsep {\bfseries #1}\hskip \labelsep {\bfseries #2.}]}{\end{trivlist}}
\newenvironment{question}[2][Question]{\begin{trivlist}
\item[\hskip \labelsep {\bfseries #1}\hskip \labelsep {\bfseries #2.}]}{\end{trivlist}}
\newenvironment{corollary}[2][Corollary]{\begin{trivlist}
\item[\hskip \labelsep {\bfseries #1}\hskip \labelsep {\bfseries #2.}]}{\end{trivlist}}
\newenvironment{solution}{\begin{proof}[Solution]}{\end{proof}}
\newenvironment{idea}[2][Proof Idea.]{\textit{#1} #2}



\parindent0in
\pagestyle{plain}
\thispagestyle{plain}

\usepackage{csquotes}
\usepackage[shortlabels]{enumitem}

\newcommand{\dated}{\today}
\newcommand{\token}[1]{\langle \text{#1} \rangle}

\begin{document}

\textbf{Introduction to the Theory of
Computation}\hfill\textbf{\myname}\\[0.01in]
\textbf{Chapter 8: Space Complexity}\hfill\textbf{\dated}\\
\smallskip\hrule\bigskip

\begin{problem}{8.29}
Show that $A_{NFA}$ is NL-complete.
\end{problem}

\begin{proof}
To show that $A_{NFA}$ is NL-complete, we must show that it is in NL and that all NL problems are log space reducible to it. To show $A_{NFA} \in$ NL, we give an algorithm $M$ that decides $A_{NFA}$ in non-deterministic log space. \\

$M =$ \textquotedblleft On input $\langle N, w \rangle$, where $N = (Q, \Sigma, \delta, q_0, F)$ is an NFA and $w$ is a string:
\begin{enumerate}
\item If some $w_i \notin \Sigma$, then \textit{reject}.
\item Let $q$, $i$ and $n$ be binary integers on the work tape, where $q$ is index of $q_0$, $n = |w|$ and $i = 0$.
\item If $i > n$, then:
\item \hspace*{0.5cm} If $q$ is an accept state in $N$, then \textit{accept}. Otherwise, \textit{reject}. 
\item Non-deterministically, repeat for binary index $q_i$ of each state in $N$:
\item \hspace*{0.5cm} If $q_i \notin \delta(q, w_i)$, then \textit{reject}.
\item \hspace*{0.5cm} Let $q = q_i$, $i = i + 1$, and go to step 2. 
\item \textit{Reject}.\textquotedblright
\end{enumerate}

To prove that $A_{NFA}$ is NL-hard, we give a log space reduction from $PATH$ to $A_{NFA}$. This reduction converts a directed graph $G$ and two nodes $s$ and $t$ to an NFA $N$ and a string $w$ so that $G$ has a directed path from $s$ to $t$, iff $N$ accepts $w$. NFA $N = (Q, \Sigma, \delta, q_0, F)$ and $w$ are constructed as follows:
\begin{enumerate}
\item Every node in $G$ is a state in $N$.
\item $\Sigma = \{1\}$.
\item Every edge $(u, v)$ in $G$ is a transition from state $u$ to $v$ by reading symbol 1. Each state in $N$ also contains a transition to it self by reading symbol 1.
\item $q_0 = s$.
\item $F = \{t\}$.
\item Finally, $w = 1^{n-1}$, where $n$ is the number of states in $N$.
\end{enumerate}

Suppose graph $G$ contains a path $P = (v_1,v_2,\dots,v_k)$ where $v_1 = s$ and $v_k = t$. Then, each edge in path $P$ is a transition in $N$. The length of string $w$ equals the maximum number of edges in a longest path in $G$. Therefore, after reading the last symbol of $w$, at least one branch of $N$ would be at accept state $t$. \\ 

Suppose NFA $N$ accepts string $w$. Then, an accepting branch of $N$ on $w$ gives a path in $G$ from $s$ to $t$.

\end{proof}
\end{document}